\documentclass[12pt,a4paper]{report}

% ============================================================
% PACKAGES
% ============================================================
\usepackage[utf8]{inputenc}
\usepackage[T1]{fontenc}
\usepackage[utf8]{vietnam}
\usepackage[T5]{fontenc}
\usepackage[vietnamese,english]{babel}
\usepackage{times}
\usepackage[margin=1in]{geometry}
\usepackage{graphicx}
\usepackage{amsmath,amssymb,amsthm}
\usepackage{algorithm}
\usepackage{algpseudocode}
\usepackage{booktabs}
\usepackage{array}
\usepackage{tabularx}
\usepackage{hyperref}
\usepackage{listings}
\usepackage{enumitem}
\usepackage{fancyhdr}
\usepackage{setspace}
\usepackage{titlesec}
\usepackage{tocloft}
\usepackage{xcolor}
\usepackage{pgfplots}
\pgfplotsset{compat=1.18}

% ============================================================
% FORMATTING
% ============================================================
\onehalfspacing

% Chapter formatting: "CHAPTER 1" then title
\titleformat{\chapter}[display]
{\normalfont\Large\bfseries\centering}
{CHAPTER \thechapter}{0pt}{\Large\uppercase}
\titlespacing*{\chapter}{0pt}{-20pt}{20pt}

% Section/subsection formatting
\titleformat{\section}{\normalfont\large\bfseries}{\thesection}{1em}{}
\titleformat{\subsection}{\normalfont\normalsize\bfseries}{\thesubsection}{1em}{}
\titleformat{\subsubsection}{\normalfont\normalsize\bfseries\itshape}{\thesubsubsection}{1em}{}

% Page numbering
\pagestyle{fancy}
\fancyhf{}
\fancyfoot[R]{\thepage}
\renewcommand{\headrulewidth}{0pt}

% TOC formatting
\renewcommand{\cftchapfont}{\bfseries}
\renewcommand{\cftsecfont}{}
\renewcommand{\cftchapleader}{\cftdotfill{\cftdotsep}}

% Hyperref setup
\hypersetup{colorlinks=true, linkcolor=black, citecolor=black, urlcolor=blue}

% Listings for pseudocode
\lstset{
  basicstyle=\small\ttfamily,
  frame=single,
  xleftmargin=1em,
  breaklines=true,
  numbers=left,
  numberstyle=\tiny,
  numbersep=5pt,
}

% Theorem environments
\newtheorem{theorem}{Theorem}[chapter]
\newtheorem{definition}{Definition}[chapter]

% ============================================================
% DOCUMENT
% ============================================================
\begin{document}

% ============================================================
% TITLE PAGE 1 (without advisor)
% ============================================================
  \thispagestyle{fancy}
  \begin{titlepage}
    \begin{center}

      VIETNAM GENERAL CONFEDERATION OF LABOUR\\
      \textbf{TON DUC THANG UNIVERSITY}\\
      \textbf{FACULTY OF INFORMATION TECHNOLOGY}\\

      \vspace{1cm}
      \includegraphics[width=4.2cm]{tdtu_logo.png}

      \vspace{2cm}
      \textbf{MÃ QUỐC CƯỜNG -- 522I0001}\\
      \textbf{NGUYỄN CAO PHI -- 521V0006}\\

      \vspace{2cm}
      {\LARGE\textbf{Object-Oriented Solutions to the Problem of Mining Top-K Maximal Frequent Itemsets from Uncertain Databases}}

      \vspace{3cm}
      {\large\textbf{Information Technology Project -- 504091}}

      \vfill
      \textbf{HO CHI MINH CITY, YEAR 2026}

    \end{center}
  \end{titlepage}

% ============================================================
% TITLE PAGE 2 (with advisor)
% ============================================================
  \thispagestyle{fancy}
  \begin{titlepage}
    \begin{center}

      VIETNAM GENERAL CONFEDERATION OF LABOUR\\
      \textbf{TON DUC THANG UNIVERSITY}\\
      \textbf{FACULTY OF INFORMATION TECHNOLOGY}\\

      \vspace{1cm}
      \includegraphics[width=4.2cm]{tdtu_logo.png}

      \vspace{2cm}
      \textbf{MÃ QUỐC CƯỜNG -- 522I0001}\\
      \textbf{NGUYỄN CAO PHI -- 521V0006}\\

      \vspace{2cm}
      {\LARGE\textbf{Object-Oriented Solutions to the Problem of Mining Top-K Maximal Frequent Itemsets from Uncertain Databases}}

      \vspace{1.5cm}
      {\large\textbf{Information Technology Project -- 504091}}

      \vspace{2cm}
      Advised by\\
      \textbf{Prof. Nguyễn Chí Thiện}

      \vfill
      \textbf{HO CHI MINH CITY, YEAR 2026}

    \end{center}
  \end{titlepage}

% ============================================================
% ACKNOWLEDGMENT
% ============================================================
  \chapter*{ACKNOWLEDGMENT}
  \addcontentsline{toc}{chapter}{ACKNOWLEDGMENT}

  We would like to express our sincere gratitude to Professor Nguyễn Chí Thiện for his patient guidance, thoughtful supervision, and continuous encouragement throughout this project. His expertise in data mining and algorithm design provided invaluable direction at every stage of our work.

  We also thank our classmates for their helpful discussions and feedback, and Ton Duc Thang University for providing the academic environment and resources that made this project possible.

  \vspace{3cm}
  \begin{flushright}
    \textit{Ho Chi Minh City, day \quad month \quad year 2026}\\
    \textit{Authors}\\[2cm]
    \textit{(Signature and full name)}
  \end{flushright}

% ============================================================
% DECLARATION OF AUTHORSHIP
% ============================================================
  \chapter*{DECLARATION OF AUTHORSHIP}
  \addcontentsline{toc}{chapter}{DECLARATION OF AUTHORSHIP}

  I hereby declare that this thesis was carried out by my team under the guidance and supervision of Prof Nguyễn Chí Thiện; and that the work and the results contained in it are original and have not been submitted anywhere for any previous purposes. The data and figures presented in this thesis are for analysis, comments, and evaluations from various resources by my own work and have been duly acknowledged in the reference part.

  In addition, other comments, reviews and data used by other authors, and organizations have been acknowledged, and explicitly cited.

  \textbf{I will take full responsibility for any fraud detected in my thesis}. Ton Duc Thang University is unrelated to any copyright infringement caused on my work (if any).

  \vspace{2cm}
  \begin{flushright}
    \textit{Ho Chi Minh City, day \quad month \quad year}\\
    \textit{Author}\\[2cm]
    \textit{[Signature and full name]}
  \end{flushright}

% ============================================================
% ABSTRACT
% ============================================================
  \chapter*{ABSTRACT}
  \addcontentsline{toc}{chapter}{ABSTRACT}

  Frequent itemset mining is a fundamental task in data mining. Classical algorithms assume binary transaction data where items are either present or absent. However, real-world databases often contain \textbf{uncertainty}: a barcode scanner might be only 80\% confident that it read a product correctly, a sensor might report a reading with known error margins, or data merged from multiple sources might contain conflicting entries. In such databases, each item is associated with an \textbf{existential probability} rather than a simple yes/no flag.

  This report addresses the problem of mining \textbf{Top-K maximal frequent itemsets} from uncertain transaction databases using object-oriented design under the expected support model. We implement and evaluate two Java algorithms:

  \textbf{UGenMax} adapts GenMax (Gouda and Zaki, 2005) using bottom-up weighted tidset intersections, and \textbf{UFPMax} adapts FPMax (Grahne and Zhu, 2003) using top-down conditional database projection. Both employ branch-and-bound search with anti-monotonicity pruning, progressive focusing, and dynamic Top-K threshold raising, and both support a Top-K mode.

  Both algorithms are implemented in pure Java SE with no external dependencies, following object-oriented design principles with separate algorithm classes (no main method) and experiment classes (with main method and statistics output). Experiments on three benchmark datasets (mushroom, retail, accidents) at medium uncertainty show that: (1) UGenMax and UFPMax produce identical results, with UGenMax being 19--43\% faster on the primary retail dataset, 42--66\% faster on mushroom, and 13--36\% faster on accidents at moderate/high thresholds; (2) both algorithms scale to 340K transactions.

  \pagebreak

% ============================================================
% TABLE OF CONTENTS
% ============================================================
  \tableofcontents
  \addcontentsline{toc}{chapter}{TABLE OF CONTENTS}

% ============================================================
% LIST OF FIGURES (placeholder)
% ============================================================
  \listoffigures
  \addcontentsline{toc}{chapter}{IMAGE TABLE}

% ============================================================
% LIST OF ABBREVIATIONS
% ============================================================
  \chapter*{List of Abbreviations}
  \addcontentsline{toc}{chapter}{List of Abbreviations}

  \begin{tabular}{@{}ll@{}}
    \toprule
    \textbf{Abbreviation} & \textbf{Definition} \\
    \midrule
    FIM & Frequent Itemset Mining \\
    MFI & Maximal Frequent Itemset \\
    DFS & Depth-First Search \\
    PWS & Possible World Semantics \\
    expSup & Expected Support \\
    minsup & Minimum Support Threshold \\
    WT & Weighted Tidset \\
    DB & Database \\
    \bottomrule
  \end{tabular}

% ============================================================
% CHAPTER 1: INTRODUCTION
% ============================================================
  \chapter{INTRODUCTION}

  \section{What is Frequent Itemset Mining?}

  Imagine a supermarket analyzing thousands of customer receipts to discover which products are frequently bought together. If many customers buy bread and butter in the same trip, the pair \{bread, butter\} is a \textbf{frequent itemset}. This information is valuable for product placement, cross-selling, and inventory management.

  More formally, given a database of transactions (where each transaction is a set of items) and a minimum support threshold, the goal of frequent itemset mining is to find all item groups that appear in at least that many transactions. This problem was first formalized by Agrawal, Imielinski, and Swami~\cite{agrawal1993} and has since become one of the most studied problems in data mining, with applications in bioinformatics, text mining, web usage analysis, network security, and recommendation systems.

  \section{The Challenge of Uncertainty}

  Classical algorithms assume perfect data: an item is either in a transaction or it is not. In practice, data is often imperfect. A temperature sensor, for example, reports values with a known error margin, so when the system logs that the sensor detected item A it might only be 85\% confident in that reading. In a data-integration setting, merging customer records from two different databases produces ambiguous matches: a merged record might indicate that a customer purchased item B, but with only 70\% confidence because of name or address mismatches. Privacy-preserving systems introduce yet another form of uncertainty by deliberately adding random noise so that each item's presence becomes probabilistic rather than certain.

  In all these cases, the result is an \textbf{uncertain transaction database} where each item in each transaction carries a probability between 0 and 1, and the challenge becomes how to mine meaningful patterns from data that we are not entirely sure about.

  \section{Why Maximal Itemsets?}

  A database with $m$ distinct items has $2^m$ possible subsets. Even a modest database with 100 items has more possible subsets than atoms in the universe. While most of these are not frequent, the number of frequent itemsets can still be enormous---potentially millions or more.

  \textbf{Maximal frequent itemsets} solve this information overload problem. A frequent itemset is maximal if no proper superset of it is also frequent. The key insight is that every frequent itemset is a subset of at least one maximal frequent itemset, so the maximal set is a \textbf{lossless compression} of the full set of frequent itemsets.

  \section{Project Objectives}

  This project has four objectives:

  \begin{enumerate}
    \item Formally define the problem of mining frequent maximal itemsets from uncertain databases under the expected support model.
    \item Design and implement two algorithms: UGenMax and UFPMax, both supporting a Top-K mode.
    \item Implement both algorithms as a pure Java SE application with no external library dependencies, with separate algorithm classes and experiment classes.
    \item Experimentally evaluate and compare the two algorithms on benchmark datasets (mushroom, retail, accidents) of varying sizes and uncertainty levels.
  \end{enumerate}

  \section{Report Structure}

  The remainder of this report is organized as follows. Chapter~2 reviews related work. Chapter~3 provides the formal problem definition for the expected support model. Chapter~4 presents both algorithms in detail. Chapter~5 describes the experimental setup. Chapter~6 presents results. Chapter~7 concludes.

% ============================================================
% CHAPTER 2: LITERATURE REVIEW
% ============================================================
  \chapter{LITERATURE REVIEW}

  \section{Frequent Itemset Mining in Deterministic Databases}

  The Apriori algorithm~\cite{agrawal1994} was the first efficient algorithm for frequent itemset mining. Its key contribution was the \textbf{anti-monotonicity property}: if an itemset is infrequent, then all of its supersets are also infrequent. The FP-Growth algorithm~\cite{han2000} improved upon Apriori by eliminating candidate generation through a compressed FP-tree data structure.

  \section{Mining Maximal Frequent Itemsets}

  Several algorithms have been developed specifically for maximal frequent itemset mining. MaxMiner~\cite{bayardo1998} was the first to use a look-ahead pruning strategy. GenMax~\cite{gouda2005} introduced depth-first backtracking with tidsets and progressive focusing, and is the basis for our UGenMax. FPMax~\cite{grahne2003} extended FP-Growth with MFI-tree superset checking, and is the basis for our UFPMax. CARPENTER-MAX~\cite{pan2003} was designed for datasets with many items but few transactions, and Charm-MFI~\cite{szathmary2006} discovers maximal itemsets by post-processing closed itemsets.

  \section{Uncertain Database Mining}

  Chui, Kao, and Hung~\cite{chui2007} pioneered uncertain data mining with the expected support model and U-Apriori algorithm. Leung et al.~\cite{leung2008} proposed UF-Growth for uncertain data. Leung and Tanbeer~\cite{leung2013} improved this with UFP-Growth. Li and Zhang~\cite{li2019} addressed the top-k variant with dynamic threshold raising. Aggarwal et al.~\cite{aggarwal2009} analyzed the theoretical complexity of uncertain pattern mining.

  \section{Probabilistic Support Mining}

  An alternative line of work defines frequency probabilistically, requiring $\Pr[\text{sup}(X) \geq T] \geq \tau$ for some integer count and probability threshold. Sun et al.~\cite{sun2010} and Chen et al.~\cite{chen2020} are representative examples; we mention them only as historical context, as the present work uses the expected support model throughout.

  \section{Gap and Our Contribution}

  Most prior work focuses on mining \textbf{all} frequent or \textbf{closed} frequent itemsets from uncertain data. Mining \textbf{maximal} frequent itemsets has received less attention, and Top-K formulations of the maximal problem under uncertainty are particularly under-explored. Our contribution fills these gaps by adapting GenMax and FPMax for uncertain data under the expected support model, equipping both with a Top-K mode and dynamic threshold raising, and providing a comprehensive experimental comparison of the two algorithms across three benchmark datasets.

% ============================================================
% CHAPTER 3: PROBLEM STATEMENT AND DEFINITIONS
% ============================================================
  \chapter{PROBLEM STATEMENT AND RELATED DEFINITIONS}

  This section provides precise mathematical definitions for each concept. We build the definitions incrementally, starting from the data model and working up to the problem statement.

  \section{Uncertain Transaction Database}

  \begin{definition}[Item Set]
Let $I = \{i_1, i_2, \ldots, i_m\}$ be a finite set of $m$ distinct items.
  \end{definition}

  \begin{definition}[Uncertain Transaction]
An uncertain transaction $t_k$ is a set of items from $I$, where each item $i$ is associated with an \textbf{existential probability} $P(i, t_k)$ satisfying:
\[ 0 < P(i, t_k) \leq 1 \]
Items are assumed to be mutually independent within each transaction.
  \end{definition}

  \begin{definition}[Uncertain Transaction Database]
An uncertain transaction database $D = \{t_1, t_2, \ldots, t_n\}$ is an ordered collection of $n$ uncertain transactions.
  \end{definition}

  \textbf{Running Example.} Throughout this report, we use the following 5-transaction database for illustration:

  \begin{table}[h]
    \centering
    \caption{The contextUncertain dataset}
    \label{tab:running_example}
    \begin{tabular}{cl}
      \toprule
      \textbf{Transaction} & \textbf{Items (with existential probabilities)} \\
      \midrule
      $t_1$ & 1(0.8) \quad 3(0.6) \quad 4(0.9) \\
      $t_2$ & 2(0.7) \quad 3(0.5) \quad 5(0.9) \\
      $t_3$ & 1(0.6) \quad 2(0.8) \quad 3(0.7) \quad 5(0.85) \\
      $t_4$ & 2(0.5) \quad 5(0.6) \\
      $t_5$ & 1(0.9) \quad 2(0.75) \quad 3(0.8) \quad 5(0.7) \\
      \bottomrule
    \end{tabular}
  \end{table}

  \section{Expected Support}

  \begin{definition}[Expected Support]
The expected support of an itemset $X = \{x_1, x_2, \ldots, x_p\}$ in database $D$ is:
\begin{equation}
  \text{expSup}(X) = \sum_{t \in D} \prod_{x \in X} P(x, t)
  \label{eq:expsup}
\end{equation}
  \end{definition}

  \textbf{Worked Example.} Consider $X = \{1, 3\}$:

  \begin{align*}
    t_1&: P(1,t_1) \times P(3,t_1) = 0.8 \times 0.6 = 0.48 \\
    t_2&: \text{item 1 absent} = 0.00 \\
    t_3&: P(1,t_3) \times P(3,t_3) = 0.6 \times 0.7 = 0.42 \\
    t_4&: \text{item 1 absent} = 0.00 \\
    t_5&: P(1,t_5) \times P(3,t_5) = 0.9 \times 0.8 = 0.72 \\
  \end{align*}
  \[ \text{expSup}(\{1, 3\}) = 0.48 + 0.00 + 0.42 + 0.00 + 0.72 = \mathbf{1.62} \]

  \section{Anti-Monotonicity Property}

  \begin{theorem}[Anti-Monotonicity of Expected Support]
For any itemset $X$ and any item $y \notin X$:
\[ \text{expSup}(X \cup \{y\}) \leq \text{expSup}(X) \]
  \end{theorem}

  \begin{proof}
    Since $0 < P(y, t) \leq 1$ for all transactions $t$, each term $\prod_{x \in X} P(x,t) \cdot P(y,t) \leq \prod_{x \in X} P(x,t)$. Therefore the sum cannot increase.
  \end{proof}

  \section{Frequent and Maximal Frequent Itemsets}

  \begin{definition}[Frequent Itemset]
Given threshold minsup, itemset $X$ is \textbf{frequent} if $\text{expSup}(X) \geq \text{minsup}$.
  \end{definition}

  \begin{definition}[Maximal Frequent Itemset]
A frequent itemset $X$ is \textbf{maximal} if for every item $y \notin X$, $\text{expSup}(X \cup \{y\}) < \text{minsup}$.
  \end{definition}

  \section{Problem Statement}

  This project addresses the following problem under the expected support model:

  \begin{itemize}
    \item \textbf{Input:} Uncertain database $D$, threshold $\text{minsup} > 0$ (or integer $K$ for Top-K mode).
    \item \textbf{Output:} All maximal frequent itemsets where $\text{expSup}(X) \geq \text{minsup}$.
    \item \textbf{Algorithms:} UGenMax, UFPMax.
  \end{itemize}

% ============================================================
% CHAPTER 4: METHODS
% ============================================================
  \chapter{METHODS}

  This chapter presents both algorithms in detail: data structures, pseudocode, worked examples, and complexity analysis.

  \section{Branch-and-Bound Framework}

  Both algorithms explore the \textbf{itemset lattice} using depth-first search (DFS). The search is a tree: the root is the empty set, each node is an itemset (prefix), and each edge adds one item. Without pruning, this requires exploring $2^m$ subsets. Branch-and-bound makes the search tractable by computing \textbf{bounds} at each node and \textbf{pruning} subtrees that cannot yield useful results.

  \subsection{Pruning 1: Anti-Monotonicity Bound}

  \textbf{Rule:} If $\text{expSup}(\text{prefix} \cup \{\text{item}\}) < \text{minsup}$, skip this item and the entire subtree below it.

  \textbf{Rationale:} By Theorem~3.1, every superset will have even lower expected support. This rule is responsible for the majority of pruning in practice. For UGenMax, a cheaper upper bound $\min(\text{expSup}(A), \text{expSup}(B))$ is checked before the full intersection.

  \subsection{Pruning 2: Progressive Focusing}

  \textbf{Rule:} Let $P$ be the current prefix and $C$ the remaining candidates. If $P \cup C \subseteq M$ for some known MFI $M$, prune the entire subtree.

  \textbf{Rationale:} Any discoverable itemset would be a subset of $M$ and therefore not maximal. This rule becomes increasingly effective as more MFIs are discovered during the search.

  \subsection{Pruning 3: Top-K Dynamic Threshold Raising}

  In Top-K mode, a \textbf{min-heap} of size $K$ tracks the best results. When a new MFI is discovered, if the heap is full and the new MFI's expected support exceeds the heap minimum, the minimum is replaced. The new heap minimum becomes the updated minsup.

  \textbf{Feedback loop:} Better results $\rightarrow$ higher threshold $\rightarrow$ more aggressive pruning $\rightarrow$ faster convergence.

  \textbf{Design difference:} UGenMax raises the threshold dynamically during search. UFPMax does \textbf{not}, because its processing order discovers single-item MFIs early, which would prematurely block multi-item MFIs.

  \section{Algorithm 1: UGenMax}

  UGenMax adapts GenMax~\cite{gouda2005} for uncertain data using \textbf{bottom-up, depth-first backtracking} with \textbf{weighted tidsets}.

  \subsection{Data Structure: Weighted Tidsets}

  \begin{definition}[Weighted Tidset]
The weighted tidset of item $i$ is:
\[ \text{WT}(i) = \{(t, P(i,t)) \mid t \in D,\ P(i,t) > 0\} \]
sorted by transaction ID. Expected support: $\text{expSup}(i) = \sum_{(t,p) \in \text{WT}(i)} p$.
  \end{definition}

  To compute $\text{expSup}(\{A, B\})$, we \textbf{intersect} $\text{WT}(A)$ and $\text{WT}(B)$ using a merge-join in $O(|\text{WT}(A)| + |\text{WT}(B)|)$ time. For each common transaction $t$, the result is $(t, P(A,t) \times P(B,t))$.

  \textbf{Worked Example.} Using items 1 and 3 from Table~\ref{tab:running_example}:

  \begin{align*}
    \text{WT}(1) &= [(t_1, 0.8),\ (t_3, 0.6),\ (t_5, 0.9)] \quad \text{expSup} = 2.3 \\
    \text{WT}(3) &= [(t_1, 0.6),\ (t_2, 0.5),\ (t_3, 0.7),\ (t_5, 0.8)] \quad \text{expSup} = 2.6
  \end{align*}

  Intersection:
  \begin{align*}
    t_1&: \text{both present} \rightarrow 0.8 \times 0.6 = 0.48 \\
    t_2&: \text{only in WT}(3) \rightarrow \text{skip} \\
    t_3&: \text{both present} \rightarrow 0.6 \times 0.7 = 0.42 \\
    t_5&: \text{both present} \rightarrow 0.9 \times 0.8 = 0.72
  \end{align*}
  \[ \text{WT}(\{1,3\}) = [(t_1, 0.48),\ (t_3, 0.42),\ (t_5, 0.72)] \quad \text{expSup} = 1.62 \]

  \subsection{Enumeration Strategy}

  Items are sorted by expected support \textbf{ascending}. Candidates are processed \textbf{most frequent first} (right to left). This ordering discovers large MFIs early, strengthening progressive focusing.

  \subsection{Pseudocode}

  \begin{algorithm}[H]
\caption{UGenMax}
\begin{algorithmic}[1]
  \Require Uncertain database $D$, threshold minsup
  \Ensure Set $\text{MFI}$ of all maximal frequent itemsets
  \State $\text{MFI} \gets \emptyset$
  \For{each item $i$ in $D$}
    \State Build $\text{WT}(i)$ by scanning all transactions
  \EndFor
  \State Remove items where $\text{expSup}(i) < \text{minsup}$
  \State Sort items by expSup ascending: $[i_1, i_2, \ldots, i_k]$
  \State \Call{GenMax}{$\emptyset$, $[i_1, \ldots, i_k]$, null}
  \State \Return MFI
\end{algorithmic}
  \end{algorithm}

  \begin{algorithm}[H]
\caption{GenMax Procedure}
\begin{algorithmic}[1]
  \Procedure{GenMax}{prefix, candidates, prefixTidset}
    \If{$\text{prefix} \cup \text{candidates} \subseteq$ some MFI in MFI} \Comment{Pruning 2}
    \State \Return
    \EndIf
    \State hasFreqExt $\gets$ false
    \For{idx $\gets$ $|$candidates$| - 1$ \textbf{downto} 0}
      \State item $\gets$ candidates[idx]
      \State newTidset $\gets$ \Call{Intersect}{prefixTidset, WT(item)}
      \State expSup $\gets$ Sum(newTidset.probabilities)
      \If{expSup $<$ minsup} \Comment{Pruning 1}
      \State \textbf{continue}
      \EndIf
      \State hasFreqExt $\gets$ true
      \State newPrefix $\gets$ prefix $\cup$ \{item\}
      \State newCands $\gets$ \{$c \in$ candidates$[0..idx{-}1]$ $|$ Intersect(newTidset, WT($c$)).expSup $\geq$ minsup\}
      \If{newCands $= \emptyset$}
        \State \Call{AddMaximal}{newPrefix, expSup}
      \Else
        \State \Call{GenMax}{newPrefix, newCands, newTidset}
      \EndIf
    \EndFor
    \If{\textbf{not} hasFreqExt \textbf{and} prefix $\neq \emptyset$}
      \State \Call{AddMaximal}{prefix, prefixTidset.expSup}
    \EndIf
  \EndProcedure
\end{algorithmic}
  \end{algorithm}

  \subsection{Complexity Analysis}

  \textbf{Time:} Worst case $O(2^m)$ nodes, each with $O(n)$ tidset intersection. Pruning dramatically reduces explored nodes in practice.

  \textbf{Space:} $O(m \times n)$ for weighted tidsets, $O(|\text{MFI}|)$ for results, $O(m)$ for the recursion stack.

  \section{Algorithm 2: UFPMax}

  UFPMax adapts FPMax~\cite{grahne2003} for uncertain data using \textbf{top-down, recursive conditional database projection}.

  \subsection{Data Structure: Projected Transactions}

  Each projected transaction is a triple (weight, items[], probs[]):
  \begin{itemize}
    \item \textbf{weight}: accumulated product of probabilities of all prefix items. Initially 1.0.
    \item \textbf{items[] and probs[]}: remaining items and probabilities, sorted by frequency-descending order.
  \end{itemize}

  Expected support of item $i$ in a conditional database:
  \begin{equation}
    \text{expSup}(i \mid \text{prefix}) = \sum_{t : i \in t} \text{weight}_t \times P(i, t)
  \end{equation}

  \subsection{Conditional Database Construction}

  When extending prefix by item $i$:
  \begin{enumerate}
    \item Select transactions containing item $i$.
    \item Compute $\text{new\_weight} = \text{old\_weight} \times P(i, t)$.
    \item Project to items \textbf{before} $i$ in header order (more frequent items).
    \item Remove items with conditional expSup $<$ minsup.
  \end{enumerate}

  \subsection{Pseudocode}

  \begin{algorithm}[H]
\caption{UFPMax}
\begin{algorithmic}[1]
  \Require Uncertain database $D$, threshold minsup
  \Ensure Set $\text{MFI}$ of all maximal frequent itemsets
  \State $\text{MFI} \gets \emptyset$
  \State Compute expSup for each single item
  \State Remove items with expSup $<$ minsup
  \State Sort items by expSup descending: header $= [h_1, \ldots, h_k]$
  \State Build initial projected DB (weight = 1.0, frequent items only)
  \State \Call{FPMax}{$\emptyset$, 0.0, projDB, header}
  \State \Return MFI
\end{algorithmic}
  \end{algorithm}

  \begin{algorithm}[H]
\caption{FPMax Procedure}
\begin{algorithmic}[1]
  \Procedure{FPMax}{prefix, prefixExpSup, condDB, headerItems}
    \If{prefix $\cup$ headerItems $\subseteq$ some MFI} \Comment{Pruning 2}
    \State \Return
    \EndIf
    \State hasFreqExt $\gets$ false
    \For{idx $\gets$ $|$headerItems$| - 1$ \textbf{downto} 0}
      \State item $\gets$ headerItems[idx]
      \State expSup $\gets \sum_{t \in \text{condDB}}$ weight$_t \times P$(item, $t$)
      \If{expSup $<$ minsup} \Comment{Pruning 1}
      \State \textbf{continue}
      \EndIf
      \State hasFreqExt $\gets$ true; newPrefix $\gets$ prefix $\cup$ \{item\}
      \State newCondDB $\gets$ \Call{BuildConditional}{condDB, item, idx}
      \State newHeader $\gets$ \{$h \in$ header$[0..idx{-}1]$ $|$ condExpSup$(h) \geq$ minsup\}
      \If{newHeader $= \emptyset$}
        \State \Call{AddMaximal}{newPrefix, expSup}
      \Else
        \State \Call{FPMax}{newPrefix, expSup, newCondDB, newHeader}
      \EndIf
    \EndFor
    \If{\textbf{not} hasFreqExt \textbf{and} prefix $\neq \emptyset$}
      \State \Call{AddMaximal}{prefix, prefixExpSup}
    \EndIf
  \EndProcedure
\end{algorithmic}
  \end{algorithm}

  \subsection{Complexity Analysis}

  \textbf{Time:} Worst case $O(2^m)$ nodes. Conditional DB construction is $O(|\text{condDB}| \times |\text{items}|)$ per node, shrinking at deeper levels.

  \textbf{Space:} $O(n \times m)$ for the initial projected DB. Each recursive level maintains its own conditional DB. Stack depth at most $m$.

  \section{Comparison of UGenMax and UFPMax}

  \begin{table}[h]
    \centering
    \caption{Comparison of UGenMax and UFPMax}
    \begin{tabularx}{\textwidth}{lXX}
\toprule
\textbf{Aspect} & \textbf{UGenMax} & \textbf{UFPMax} \\
\midrule
Base algorithm & GenMax (Gouda \& Zaki, 2005) & FPMax (Grahne \& Zhu, 2003) \\
Approach & Bottom-up (vertical) & Top-down (horizontal) \\
Data structure & Weighted tidsets & Projected conditional DBs \\
Item ordering & Ascending expSup & Descending expSup \\
Top-K threshold & Raised dynamically & Fixed, select after search \\
Best for & Sparse data & Dense data \\
\bottomrule
    \end{tabularx}
  \end{table}

  \section{Top-K Seeding}

  Top-K mode is conceptually simple --- ``return the $K$ itemsets with the highest expected support'' --- but a naive implementation collapses for a subtle reason: with no minsup threshold, the initial value of \texttt{minsup} is $0$, every single item is frequent, and the anti-monotonicity pruning rule (Pruning~1) cannot fire. The search degenerates into an unpruned enumeration of the full lattice, which is exactly what the branch-and-bound framework was designed to avoid. The Top-K seeding phase exists to break this deadlock by establishing a non-trivial initial threshold \emph{before} the main search begins.

  \paragraph{The deadlock and the fix.} The dynamic threshold-raising rule (Pruning~3) only kicks in once the Top-K heap is full, because the heap minimum is what gets fed back into \texttt{minsup}. The cheapest way to fill the heap is to pre-compute scores for a small set of itemsets that we already know are cheap to evaluate, namely single items and 2-item pairs over a bounded set of candidates. These items can be ranked directly from the item-level expected supports and from a single tidset intersection per pair --- no recursion required. Once these scores are in the heap, its minimum becomes a usable threshold and Pruning~3 starts firing on the first depth-first step.

  \paragraph{Procedure.} Both UGenMax and UFPMax run the same five-step seeding phase before entering their respective recursive enumeration:

  \begin{enumerate}
    \item \textbf{Rank items.} Sort every distinct item by expected support in descending order and keep the top $50$ most frequent. The cap of $50$ is deliberate --- the pairing step is $O(50^2) = 2{,}500$ intersections in the worst case, which is negligible compared to the main search, while extending the cap to all items would make seeding itself $O(m^2)$ on the full item universe.
    \item \textbf{Enumerate cheap candidates.} For each of the $50$ items, push the singleton onto a candidate list. For each pair $(i,j)$ with $i<j$, compute the intersection of their weighted tidsets and push the resulting 2-itemset onto the same list if its expected support is positive. After this step the candidate list contains up to $50 + \binom{50}{2} = 1{,}275$ itemsets, each with a known score.
    \item \textbf{Fill the Top-K heap.} Sort the candidate list by expected support descending and insert the highest-scoring ones into a fresh \texttt{TopKHeap} of capacity $K$, skipping duplicates. Once the heap holds $K$ entries, its minimum is the lowest score we are willing to keep at this point in the run.
    \item \textbf{Seed minsup at 50\% of the heap minimum.} Set the working \texttt{minsup} to half of the heap minimum, rather than the heap minimum itself. The reason is that the seed itemsets are restricted to size $\leq 2$, but the eventual top-$K$ MFIs may have many more items and therefore lower expected support (anti-monotonicity). Using the raw heap minimum as \texttt{minsup} would prune away the very supersets we are searching for; the $50\%$ slack is a heuristic margin that empirically retains enough of the search space to find the true top-$K$ on every dataset we tested.
    \item \textbf{Reset the heap and start the main search.} Discard the seed heap and create a fresh one of capacity $K$. The seeded itemsets were only used to set \texttt{minsup}; they are not themselves treated as MFIs, because we have not yet checked them for maximality and because their scores were computed during the seeding phase without any guarantee that they survive the maximality filter. The main DFS then runs as usual, with Pruning~3 immediately active.
  \end{enumerate}

  \paragraph{Retry loop.} The $50\%$ slack is intentionally conservative, but it can occasionally overshoot. If the main search finishes with fewer than $K$ MFIs in the final heap --- typically because the seed threshold was tighter than the true top-$K$ allowed --- the algorithm halves \texttt{minsup} \emph{again} (i.e.\ uses $25\%$ of the previous base) and re-runs the main search from scratch with a fresh heap. This is repeated up to $3$ times, giving effective thresholds of $0.5\times$, $0.125\times$, $0.0625\times$, $0.0156\times$ relative to the seed. In our experiments the first attempt succeeds on every (\texttt{dataset}, $K$) pair tested, but the retry loop is what guarantees termination with $|\text{result}| = K$ whenever at least $K$ MFIs exist.

  \paragraph{Cost.} Seeding adds $O(s + s^2 \cdot n)$ work, where $s = 50$ is the seed cap and $n$ is the average tidset length. On all three benchmark datasets this takes well under one second, while the gain from having Pruning~3 active during the main DFS is several orders of magnitude. The combined effect is what makes Top-K mode tractable at all --- without seeding, even mushroom (the smallest dataset at $8{,}416$ transactions) would not finish a Top-$5$ query in the time budget.

  \section{Maximality Checking}

  When a leaf node is reached, the candidate MFI is checked: (1) it must not be a subset of any existing MFI, and (2) existing MFIs that are subsets of the new candidate are removed. Both checks use sorted arrays for $O(|A| + |B|)$ merge-based comparison.

% ============================================================
% CHAPTER 5: EXPERIMENT DESIGN
% ============================================================
  \chapter{EXPERIMENT DESIGN}

  \section{Implementation Details}

  Both algorithms are implemented in Java SE 25 with no external dependencies. The code is organized into packages:
  \begin{itemize}
    \item \texttt{algorithm/}: UGenMax, UFPMax, WeightedTidset
    \item \texttt{model/}: UncertainDatabase, UncertainTransaction, Itemset, TopKHeap
    \item \texttt{util/}: PerformanceTracker
    \item \texttt{experiment/}: ExperimentRunner, ExpUFPMax, ExpUGenMax
  \end{itemize}

  Algorithm classes contain only the algorithm logic (no \texttt{main} method). Experiment classes contain \texttt{main} methods that load the database, run the algorithm, and output results with performance statistics. Every class includes Javadoc documentation.

  \section{Datasets}

  \begin{table}[h]
    \centering
    \caption{Experimental datasets}
    \begin{tabular}{lcccl}
      \toprule
      \textbf{Dataset} & \textbf{Trans.} & \textbf{Items} & \textbf{Density} & \textbf{Source} \\
      \midrule
      contextUncertain & 5 & 5 & Dense & Manual \\
      mushroom & 8,416 & 120 & Dense & UCI \\
      retail & 88,162 & 16,470 & Sparse & SPMF \\
      accidents & 340,183 & 468 & Very dense & SPMF \\
      \bottomrule
    \end{tabular}
  \end{table}

  Each standard dataset is converted to uncertain format at medium uncertainty using the \texttt{DatasetConverter} utility: each item probability is drawn from a truncated normal distribution $N(\mu{=}0.5, \sigma{=}0.2)$ clamped to $[0.1, 1.0]$. This simulates moderately noisy/uncertain transaction data and follows standard practice in uncertain data mining research~\cite{chui2007,leung2008}. All standard datasets are publicly available from the SPMF data mining library.\footnote{\url{https://www.philippe-fournier-viger.com/spmf/index.php?link=datasets.php}}

  \section{Configurations}

  \textbf{Static minsup mode (UGenMax, UFPMax).} Vary minsup across an appropriate range per dataset --- a coarse grid of six values chosen so that the MFI count varies from hundreds to single digits. Mushroom uses minsup $\in \{500, 1000, 1500, 2000, 2500, 3000\}$, retail uses $\{1000, 2000, 3000, 5000, 8000, 10000\}$, and accidents uses $\{20000, 30000, 40000, 50000, 60000, 70000\}$. The full grid is recorded in \texttt{results/exp1\_expected\_support.csv}.

  \textbf{Top-K mode (UGenMax, UFPMax).} Vary $K \in \{10, 20, 50\}$ on mushroom\_medium, recorded in \texttt{results/exp\_topk.csv} (the \texttt{ExperimentRunner} class also supports $K \in \{5, 100\}$ and the larger retail/accidents datasets, but those runs are out of scope for the headline measurements reported in Chapter~6).

  \textbf{Cross-dataset snapshot.} A single representative minsup per dataset (mushroom=2{,}000, retail=5{,}000, accidents=50{,}000) is re-run across all three datasets to compare scalability at one fixed working point, recorded in \texttt{results/exp4\_scalability.csv}.

  For each configuration we record: time (ms), memory (MB), number of MFIs, and algorithm-specific statistics (nodes explored, pruning counts, candidates generated --- emitted by the per-algorithm \texttt{ExpUFPMax} and \texttt{ExpUGenMax} runners, which is where the ``Nodes Explored'' column in Table~\ref{tab:topk} originates).

  \section{Correctness Validation}

  Correctness is validated in two ways. Both algorithms are first hand-verified on the small contextUncertain example. UGenMax and UFPMax are then cross-checked against each other on every configuration to confirm they produce identical MFI sets.

  \section{Performance Measurement}

  PerformanceTracker uses \texttt{System.nanoTime()} for timing and \texttt{Runtime.totalMemory() - Runtime.freeMemory()} for memory, with GC forced before each measured run. The full benchmark suite (\texttt{ExperimentRunner}) averages each (dataset, algorithm, parameter) cell over 3 measured runs after 1 warm-up run, for both the static-minsup grid and the Top-K grid. The cross-dataset and uncertainty-level snapshots are single-run measurements --- their purpose is to show qualitative scaling, not to claim tight error bars. All CSVs emitted by the runner live under \texttt{results/} and are the sole source for every table and chart in Chapter~6.

% ============================================================
% CHAPTER 6: EXPERIMENT RESULTS
% ============================================================
  \chapter{EXPERIMENT RESULTS}

  \section{Correctness Verification}

  On contextUncertain, both algorithms produce correct results verified by hand computation. UGenMax and UFPMax produce \textbf{identical MFI sets} across all datasets and all minsup values tested, confirming correctness.

  \section{Expected Support on retail\_medium (Primary Dataset)}

  Retail is our primary benchmark for the expected-support model: 88,162 transactions over 16,470 distinct items, with an average transaction length of $\approx 10$ items. It is the largest sparse dataset among our benchmarks and exercises the sparse-data regime that industrial market-basket applications most often encounter. Table~\ref{tab:retail_minsup} shows UFPMax and UGenMax on retail\_medium across a wide range of minsup thresholds.

  \begin{table}[h]
    \centering
    \caption{Performance on retail\_medium (88,162 transactions, 16,470 items)}
    \label{tab:retail_minsup}
    \begin{tabular}{rrrrrr}
      \toprule
      \textbf{minsup} & \textbf{MFIs} & \multicolumn{2}{c}{\textbf{Time (ms)}} & \multicolumn{2}{c}{\textbf{Memory (MB)}} \\
      & & UFPMax & UGenMax & UFPMax & UGenMax \\
      \midrule
      1000 & 18 & 164 & 93 & 90.3 & 67.0 \\
      2000 & 6 & 88 & 71 & 102.0 & 64.0 \\
      3000 & 4 & 82 & 60 & 100.0 & 62.0 \\
      5000 & 4 & 86 & 65 & 100.0 & 62.0 \\
      8000 & 2 & 75 & 58 & 91.0 & 59.0 \\
      10000 & 2 & 72 & 56 & 91.0 & 59.0 \\
      \bottomrule
    \end{tabular}
  \end{table}

  \textbf{Observations on retail\_medium.} As minsup rises, both the MFI count and execution time decrease, consistent with anti-monotonicity pruning. UGenMax is faster than UFPMax at every minsup tested, with speedups ranging from 19\% at minsup=2000 (71 vs.\ 88 ms) up to 43\% at minsup=1000 (93 vs.\ 164 ms), and it uses roughly 30--35\% less memory because the weighted-tidset representation is more compact than UFPMax's projected conditional databases on this sparse dataset. Absolute runtime stays well under 200 ms even at the most permissive threshold, because most of the 16,470 items are infrequent and pruned early --- the cost is dominated by database loading and tidset/projection construction rather than by the search itself. At high minsup ($\geq 3000$) only 2--4 MFIs are found, which is typical for a sparse transactional dataset where long co-occurrence patterns are rare.

  \begin{figure}[h]
    \centering
    \begin{tikzpicture}
      \begin{axis}[
        width=0.85\textwidth,
        height=6.5cm,
        xlabel={Minimum Expected Support (minsup)},
        ylabel={Execution Time (ms)},
        legend pos=north east,
        legend style={font=\small},
        grid=major,
        grid style={dashed, gray!30},
        ymin=0, ymax=200,
        xtick={1000,2000,3000,5000,8000,10000},
        xticklabel style={font=\small},
        yticklabel style={font=\small},
        mark size=2.5pt,
        thick,
      ]
        \addplot[color=blue, mark=square*] coordinates {
          (1000,164) (2000,88) (3000,82) (5000,86) (8000,75) (10000,72)
        };
        \addplot[color=red, mark=triangle*] coordinates {
          (1000,93) (2000,71) (3000,60) (5000,65) (8000,58) (10000,56)
        };
        \legend{UFPMax, UGenMax}
      \end{axis}
    \end{tikzpicture}
    \caption{Execution time vs.\ minsup on retail\_medium (88,162 transactions)}
    \label{fig:retail_time}
  \end{figure}

  \section{Expected Support on mushroom\_medium (Secondary Dataset)}

  To confirm that the trends observed on retail generalise beyond the sparse regime, we repeat the experiment on the dense mushroom dataset (Table~\ref{tab:mushroom_minsup}).

  \begin{table}[h]
    \centering
    \caption{Performance on mushroom\_medium (8,416 transactions, 120 items)}
    \label{tab:mushroom_minsup}
    \begin{tabular}{rrrrrr}
      \toprule
      \textbf{minsup} & \textbf{MFIs} & \multicolumn{2}{c}{\textbf{Time (ms)}} & \multicolumn{2}{c}{\textbf{Memory (MB)}} \\
      & & UFPMax & UGenMax & UFPMax & UGenMax \\
      \midrule
      500 & 314 & 236 & 134 & 45.5 & 14.7 \\
      1000 & 73 & 86 & 49 & 6.0 & 17.0 \\
      1500 & 30 & 53 & 25 & 18.7 & 11.1 \\
      2000 & 14 & 29 & 10 & 9.2 & 19.1 \\
      2500 & 9 & 16 & 7 & 27.4 & 12.5 \\
      3000 & 5 & 12 & 7 & 18.7 & 9.6 \\
      \bottomrule
    \end{tabular}
  \end{table}

  \textbf{Observations on mushroom\_medium:} The same qualitative trend holds --- UGenMax is consistently faster (42--66\% speedup) across all minsup values, and memory use is also lower or comparable. Because mushroom is much denser (120 items, $\sim$23 items per transaction), it produces far more MFIs at comparable relative thresholds (up to 314 MFIs at minsup=500), giving a larger absolute speedup window than the sparser retail dataset.

  \begin{figure}[h]
    \centering
    \begin{tikzpicture}
      \begin{axis}[
        width=0.85\textwidth,
        height=6.5cm,
        xlabel={Minimum Expected Support (minsup)},
        ylabel={Execution Time (ms)},
        legend pos=north east,
        legend style={font=\small},
        grid=major,
        grid style={dashed, gray!30},
        ymin=0,
        xtick={500,1000,1500,2000,2500,3000},
        xticklabel style={font=\small},
        yticklabel style={font=\small},
        mark size=2.5pt,
        thick,
      ]
        \addplot[color=blue, mark=square*] coordinates {
          (500,236) (1000,86) (1500,53) (2000,29) (2500,16) (3000,12)
        };
        \addplot[color=red, mark=triangle*] coordinates {
          (500,134) (1000,49) (1500,25) (2000,10) (2500,7) (3000,7)
        };
        \legend{UFPMax, UGenMax}
      \end{axis}
    \end{tikzpicture}
    \caption{Execution time vs.\ minsup on mushroom\_medium (8,416 transactions)}
    \label{fig:mushroom_time}
  \end{figure}

  \section{Scalability: Accidents Dataset}

  Table~\ref{tab:accidents} shows results on the accidents dataset, which is the largest and densest.

  \begin{table}[h]
    \centering
    \caption{Performance on accidents\_medium (340,183 transactions, 468 items)}
    \label{tab:accidents}
    \begin{tabular}{rrrrrr}
      \toprule
      \textbf{minsup} & \textbf{MFIs} & \multicolumn{2}{c}{\textbf{Time (s)}} & \multicolumn{2}{c}{\textbf{Memory (MB)}} \\
      & & UFPMax & UGenMax & UFPMax & UGenMax \\
      \midrule
      20,000 & 1,285 & 47.4 & 48.8 & 1,400.7 & 1,023.8 \\
      30,000 & 453 & 23.2 & 20.2 & 991.9 & 1,060.9 \\
      40,000 & 236 & 14.1 & 10.7 & 803.0 & 1,118.3 \\
      50,000 & 154 & 9.1 & 6.6 & 1,105.0 & 959.5 \\
      60,000 & 87 & 6.1 & 4.2 & 917.3 & 988.8 \\
      70,000 & 48 & 4.2 & 2.7 & 872.0 & 1,095.3 \\
      \bottomrule
    \end{tabular}
  \end{table}

  \textbf{Observations.} Both algorithms scale to 340K transactions, returning results in seconds to tens of seconds. UGenMax is 13--36\% faster than UFPMax at minsup $\geq$ 30,000 on this very dense dataset; at the lowest threshold (minsup=20,000) the two algorithms are essentially tied (UGenMax is marginally slower, 48.8 s vs.\ 47.4 s), because the very large result set of 1,285 MFIs erases most of the progressive-focusing advantage. That the search at minsup=20,000 still completes in under a minute while producing 1,285 MFIs shows that both algorithms handle large result sets effectively. Memory usage reaches $\sim$1.4 GB at low thresholds due to the large number of items per transaction ($\sim$33 items on average out of 468); here UFPMax and UGenMax use comparable memory (803--1,401 MB and 960--1,118 MB respectively).

  \begin{figure}[h]
    \centering
    \begin{tikzpicture}
      \begin{axis}[
        width=0.85\textwidth,
        height=6.5cm,
        xlabel={Minimum Expected Support (minsup)},
        ylabel={Execution Time (seconds)},
        legend pos=north east,
        legend style={font=\small},
        grid=major,
        grid style={dashed, gray!30},
        ymin=0,
        xtick={20000,30000,40000,50000,60000,70000},
        xticklabel style={font=\small, /pgf/number format/1000 sep={\,}},
        yticklabel style={font=\small},
        mark size=2.5pt,
        thick,
      ]
        \addplot[color=blue, mark=square*] coordinates {
          (20000,47.4) (30000,23.2) (40000,14.1) (50000,9.1) (60000,6.1) (70000,4.2)
        };
        \addplot[color=red, mark=triangle*] coordinates {
          (20000,48.8) (30000,20.2) (40000,10.7) (50000,6.6) (60000,4.2) (70000,2.7)
        };
        \legend{UFPMax, UGenMax}
      \end{axis}
    \end{tikzpicture}
    \caption{Execution time vs.\ minsup on accidents\_medium (340,183 transactions)}
    \label{fig:accidents_time}
  \end{figure}

  \section{Cross-Dataset Scalability Snapshot}

  Table~\ref{tab:scalability} summarises a single-run comparison across all three standard datasets at a representative minsup per dataset, highlighting how the two expected-support algorithms behave as both the transaction count and item count grow.

  \begin{table}[h]
    \centering
    \caption{Cross-dataset scalability of UFPMax and UGenMax at medium uncertainty}
    \label{tab:scalability}
    \begin{tabular}{lrrrrrrr}
      \toprule
      \textbf{Dataset} & \textbf{\#Trans} & \textbf{\#Items} & \textbf{minsup} & \textbf{MFIs} & \textbf{Algo} & \textbf{Time (ms)} & \textbf{Mem (MB)} \\
      \midrule
      mushroom & 8,416 & 120 & 2,000 & 14 & UFPMax & 41 & 22.1 \\
      mushroom & 8,416 & 120 & 2,000 & 14 & UGenMax & 14 & 23.9 \\
      retail & 88,162 & 16,470 & 5,000 & 4 & UFPMax & 127 & 130.0 \\
      retail & 88,162 & 16,470 & 5,000 & 4 & UGenMax & 64 & 82.0 \\
      accidents & 340,183 & 468 & 50,000 & 154 & UFPMax & 11,916 & 1,552.0 \\
      accidents & 340,183 & 468 & 50,000 & 154 & UGenMax & 8,222 & 1,180.0 \\
      \bottomrule
    \end{tabular}
  \end{table}

  \textbf{Observation:} Both algorithms scale from 8K to 340K transactions, and UGenMax retains a clear advantage at every scale (66\% faster on mushroom, 50\% faster on retail, 31\% faster on accidents in this snapshot). The retail run in particular illustrates UGenMax's memory advantage on sparse data ($\sim$37\% less memory than UFPMax).

  \begin{figure}[h]
    \centering
    \begin{tikzpicture}
      \begin{axis}[
        width=0.85\textwidth,
        height=6.5cm,
        ybar,
        bar width=20pt,
        xlabel={Dataset},
        ylabel={Execution Time (ms, log scale)},
        legend pos=north west,
        legend style={font=\small},
        grid=major,
        grid style={dashed, gray!30},
        ymode=log,
        log ticks with fixed point,
        ymin=5,
        symbolic x coords={mushroom, retail, accidents},
        xtick=data,
        xticklabel style={font=\small},
        yticklabel style={font=\small},
        nodes near coords,
        nodes near coords style={font=\tiny},
        every node near coord/.append style={yshift=3pt},
      ]
        \addplot[fill=blue!60] coordinates {
          (mushroom, 41)
          (retail, 127)
          (accidents, 11916)
        };
        \addplot[fill=red!60] coordinates {
          (mushroom, 14)
          (retail, 64)
          (accidents, 8222)
        };
        \legend{UFPMax, UGenMax}
      \end{axis}
    \end{tikzpicture}
    \caption{Cross-dataset scalability: UFPMax vs.\ UGenMax (log scale)}
    \label{fig:scalability}
  \end{figure}

  \section{Top-K Mode}

Top-K mode is the practitioner-friendly setting: instead of guessing a minsup threshold, the user simply asks for the $K$ highest-support MFIs. Table~\ref{tab:topk} compares UGenMax and UFPMax on mushroom\_medium across $K \in \{10, 20, 50\}$. The time and memory columns are averaged over the same 3-run measurement protocol as the static-minsup experiments; the ``Nodes Explored'' and ``Final minsup'' columns come from the per-algorithm \texttt{ExpUFPMax}/\texttt{ExpUGenMax} runners (which emit those internal counters) rather than from the bulk CSV.

\begin{table}[h]
\centering
\caption{Top-K mode on mushroom\_medium (8,416 transactions, 120 items)}
\label{tab:topk}
\begin{tabular}{rrrrrrrr}
\toprule
\textbf{K} & \multicolumn{2}{c}{\textbf{Time (ms)}} & \multicolumn{2}{c}{\textbf{Memory (MB)}} & \multicolumn{2}{c}{\textbf{Nodes Explored}} & \textbf{Final minsup} \\
 & UFPMax & UGenMax & UFPMax & UGenMax & UFPMax & UGenMax & UGenMax \\
\midrule
10 & 700 & 139 & 51.7 & 17.0 & 10 & 5 & 1,975.7 \\
20 & 725 & 144 & 24.7 & 43.6 & 17 & 6 & 1,332.9 \\
50 & 937 & 191 & 64.5 & 44.1 & 44 & 15 & 955.1 \\
\bottomrule
\end{tabular}
\end{table}

\textbf{Key findings:}
\begin{itemize}
  \item UGenMax is \textbf{4--5$\times$ faster} than UFPMax in Top-K mode across all values of $K$, because its dynamic threshold raising (Pruning~3) aggressively increases minsup during search --- for example at $K{=}10$ the threshold rises from the seeded 1,226 to a final 1,976, pruning subtrees that UFPMax must still explore.
  \item UGenMax explores far fewer nodes (5--15 vs.\ 10--44), confirming that dynamic threshold raising is highly effective.
  \item UFPMax does not raise minsup during search (its final threshold stays at the seeded value), so it discovers many more candidate MFIs (41--192) before selecting the Top-K, whereas UGenMax's rising threshold limits candidates to 22--102.
  \item As $K$ increases, both algorithms take longer because the threshold stabilises at a lower value, expanding the search space.
\end{itemize}

  \section{Summary of Findings}

  Taking the Chapter~6 results together, three findings stand out. First, correctness holds throughout: UGenMax and UFPMax produce identical MFI sets across all datasets and thresholds. Second, on the primary retail\_medium dataset UGenMax is 19--43\% faster than UFPMax and uses roughly 30--35\% less memory --- an advantage that generalises to dense mushroom (42--66\% faster) and to very dense accidents (13--36\% faster at minsup $\geq$ 30,000, essentially tied at minsup=20,000). Finally, both algorithms scale to 340K transactions (accidents) within acceptable time limits, and UGenMax additionally outperforms UFPMax in Top-K mode (4--5$\times$ faster) thanks to dynamic threshold raising.

% ============================================================
% CHAPTER 7: CONCLUSIONS
% ============================================================
  \chapter{CONCLUSIONS}

  \section{Summary}

  This project designed, implemented, and evaluated two algorithms for mining frequent maximal itemsets from uncertain transaction databases under the expected support model. UGenMax (bottom-up weighted tidset intersections, adapted from GenMax) and UFPMax (top-down conditional database projection, adapted from FPMax) both use branch-and-bound with three pruning strategies: anti-monotonicity, progressive focusing, and dynamic Top-K threshold raising. Both algorithms also support a Top-K mode.

  Experiments on three benchmark datasets (mushroom, retail, accidents) at medium uncertainty, with retail\_medium as the primary benchmark, demonstrate that both algorithms produce correct results and exhibit distinct performance characteristics.

  \section{Key Contributions}

  The project's main contributions are as follows. We formalised the problem with precise mathematical definitions for the expected support model. We adapted GenMax and FPMax for uncertain data, equipping both with a Top-K mode and dynamic threshold raising. The whole code base is pure Java SE with no external dependencies, organised into algorithm classes (no \texttt{main}) and experiment classes (with \texttt{main} and statistics output). Finally, we carried out a comprehensive experimental comparison across three datasets, two algorithms, and varying parameter configurations at medium uncertainty, showing that UGenMax consistently outperforms UFPMax in both static-minsup and Top-K modes.

  \section{Limitations}

  Several limitations remain. Maximality checking uses a linear scan ($O(|\text{MFI}|)$ per check), which a trie- or hash-based index could improve. Memory usage on the accidents dataset also reaches roughly 1.4\,GB at low thresholds, which may be restrictive in memory-constrained environments.

  \section{Future Work}

  A natural next step is a diffset optimisation for UGenMax to reduce memory usage, followed by parallel search across independent subtrees for multi-core speedup. The framework could also be extended to mining closed or sequential patterns from uncertain data, and evaluated on real-world uncertain datasets from sensor networks or noisy data integration.

% ============================================================
% REFERENCES
% ============================================================
  \begin{thebibliography}{99}

    \bibitem{agrawal1993}
    Agrawal, R., Imielinski, T., and Swami, A. (1993). Mining association rules between sets of items in large databases. \textit{Proc. ACM SIGMOD}, pp.~207--216.

    \bibitem{agrawal1994}
    Agrawal, R. and Srikant, R. (1994). Fast algorithms for mining association rules. \textit{Proc. VLDB}, pp.~487--499.

    \bibitem{aggarwal2009}
    Aggarwal, C.C., Li, Y., Wang, J., and Wang, J. (2009). Frequent pattern mining with uncertain data. \textit{Proc. ACM SIGKDD}, pp.~29--38.

    \bibitem{bayardo1998}
    Bayardo, R.J. (1998). Efficiently mining long patterns from databases. \textit{Proc. ACM SIGMOD}, pp.~85--93.

    \bibitem{chui2007}
    Chui, C.K., Kao, B., and Hung, E. (2007). Mining frequent itemsets from uncertain data. \textit{Proc. PAKDD}, pp.~47--58.

    \bibitem{fournier2014}
    Fournier-Viger, P., et al. (2014). SPMF: A Java open-source pattern mining library. \textit{JMLR}, 15, pp.~3569--3573.

    \bibitem{gouda2005}
    Gouda, K. and Zaki, M.J. (2005). GenMax: An efficient algorithm for mining maximal frequent itemsets. \textit{Data Mining and Knowledge Discovery}, 11(3), pp.~223--242.

    \bibitem{grahne2003}
    Grahne, G. and Zhu, J. (2003). Efficiently using prefix-trees in mining frequent itemsets. \textit{Proc. IEEE ICDM Workshop on FIMI}.

    \bibitem{han2000}
    Han, J., Pei, J., and Yin, Y. (2000). Mining frequent patterns without candidate generation. \textit{Proc. ACM SIGMOD}, pp.~1--12.

    \bibitem{leung2008}
    Leung, C.K.-S., Mateo, M.A.F., and Brajczuk, D.A. (2008). A tree-based approach for frequent pattern mining from uncertain data. \textit{Proc. PAKDD}, pp.~653--661.

    \bibitem{leung2013}
    Leung, C.K.-S. and Tanbeer, S.K. (2013). Fast tree-based mining of frequent itemsets from uncertain data. \textit{Proc. DASFAA}, pp.~272--287.

    \bibitem{li2019}
    Li, H. and Zhang, N. (2019). Mining top-k frequent itemsets from uncertain databases. \textit{Proc. DSAA}.

    \bibitem{pan2003}
    Pan, F., Cong, G., Tung, A.K.H., Yang, J., and Zaki, M.J. (2003). CARPENTER: Finding closed patterns in long biological datasets. \textit{Proc. ACM SIGKDD}, pp.~637--642.

    \bibitem{sun2010}
    Sun, L., Cheng, R., Cheung, D.W., and Cheng, J. (2010). Mining uncertain data with probabilistic guarantees. \textit{Proc. ACM SIGKDD}, pp.~273--282.

    \bibitem{chen2020}
    Chen, S., Nie, L., Tao, X., Li, Z., and Zhao, L. (2020). Approximation of probabilistic maximal frequent itemset mining over uncertain sensed data. \textit{IEEE Access}, 8, pp.~97529--97539.

    \bibitem{szathmary2006}
    Szathmary, L., Napoli, A., and Valtchev, P. (2006). Towards rare itemset mining. \textit{Proc. IEEE ICTAI}.

  \end{thebibliography}

\end{document}